\scanpage{36}
\mspage{27}
% archive sheet number: 36
% inv: text | top | En effet b) implique
En effet, b) implique déjà que $\varphi(I)$ est fermée. Il reste à prouver que \deleted{\unsure{les} \unsure{figures} $X_i$} pour $i,j\in I$, on a $X_i\unsure{\gtrless|\bowtie}X_j$. Donc que
% inv: display | upper third | X prime in tilde Xi setminus
\[
X'\in\widetilde X_i\setminus\widetilde X_i\cap\widetilde X_j,\quad Y'\in\widetilde X_j\setminus\widetilde X_i\cap\widetilde X_j\Longrightarrow X'\mathrel{|\circ|}Y'.
\]
% inv: text | upper third | Or par b)
Or par \deleted{\unsure{a)}} b) on aura $X'=X_{i'}$, avec $i'\leq i$, idem pour $Y'=\deleted{\unsure{Y}}X_{j'}$, avec $j'\leq j$ ; d'ailleurs comme $X'\ne Y'$ on aura $i'\ne j'$ par a), Donc on a $X_{i'}\mathrel{|\circ|}X_{j'}$ par c), qed.

% inv: text | middle | Variante polyédrale
\deleted{Morphismes de \unsure{mag}} \inserted{Variante polyédrale :} On suppose maintenant l'ens. ordonné $I$ polyédral. On veut se borner aux figures polyédrales. Cela \deleted{\unsure{est une}} \unsure{se fera} sans frais ; \deleted{\unsure{tout}} \unsure{ceci est|sera} automatique, puisque la condition polyédrale sur une figure se lit sur sa maquette \deleted{$\unsure{\widetilde E}$} $\widetilde F$, et signifie justement que celle-ci est un ens. ordonné polyédral.

% inv: text | lower middle | Morphismes de maquettes
\uline{Morphismes de maquettes}\quad $\varphi:M\longrightarrow M'$ : application telle que si $X,Y\in M$, $X'=\varphi(X)$, $Y'=\varphi(Y)$, on ait
% inv: display | lower third | a) b) c) morphism conditions
\[
\begin{aligned}
a)\quad X\leq Y&\Longrightarrow X'\leq Y'\\
b)\quad X\overset{\circ}{\ll}Y&\Longrightarrow X'\overset{\circ}{\ll}Y'\\
c)\quad X\mathrel{|\circ|}Y&\Longrightarrow X'\mathrel{|\circ|}Y'
\end{aligned}
\]
% inv: insertion | after c) | plus d) plus loin page suiv
\inserted{[+ \textcircled{d} plus loin page suiv.]}
% inv: margin note | right beside a) b) | varphi croissante
\margin{$\varphi$ croissante pour $\leq$, $\overset{\circ}{\ll}$ (alors aussi pour $\ll$)}
% inv: text | lower third | on en conclut
on en conclut
% inv: display | lower third | X ll Y implies
\[
X\ll Y\Longrightarrow X'\ll Y'
\]
% inv: text | bottom | Mais aura-t-on aussi
Mais aura-t-on aussi
% inv: display | bottom | questioned implications
\[
\begin{aligned}
X\gtrless Y&\overset{?}{\Longrightarrow}X'\gtrless Y'\\
X\parallel Y&\overset{?}{\Longrightarrow}X'\gtrless Y'
\end{aligned}
\]
% inv: margin note | slanted left lower | Peut-être vaut-il mieux
% margin note written sideways
% slanted
\margin{Peut-être vaut-il mieux parler de \unsure{comorphismes} ? Ça \unsure{ressemble} plus aux propriétés d'une application ``\unsure{image|images} \unsure{inverse|inverses}'' ! Mais alors, ne serait-il pas plus naturel de la regarder tout de suite comme $\widetilde F\longrightarrow\widetilde F'$ des \unsure{étalages|étalées}, sans supposer \unsure{fixé un} $M\longrightarrow M'$ ??}
% inv: margin note | right lower | Est-on trop exigeant
\margin{(Est-on trop exigeant ??) Ça semble raisonnable dans le cas d'un \uuline{plongement} des maquettes \unsure{surtout|seulement} $\rightarrow$. On peut songer à demander \deleted{\unsure{aussi}} \inserted{seulement} $X\gtrless Y\deleted{\Longleftarrow}\Longrightarrow X'\gtrless Y'$ ce qui est moins extravagant\ldots}
% inv: margin note | bottom right | NB Si M est un magasin trivial
\margin{\uline{NB} Si $M$ est un magasin trivial, ces conditions impliquent $\varphi$ injective.}
